\section{記法と量子力学の公理}

\noindent\textbf{この章の中心命題}
\[
\boxed{
\psi\in\mathcal{H},\qquad
\hat{A}=\hat{A}^\dagger,\qquad
\|\psi\|^2=\langle\psi,\psi\rangle=1
}
\]
量子力学では、準備された状態をヒルベルト空間のベクトルで表し、測定で問う物理量をその空間上の自己共役演算子で表す。

前節の破綻は、少なくとも次の三点を古典力学とは別の仕方で定め直すことを要求している。すなわち、状態を何で表すか、物理量を何で表すか、そして時間発展をどう与えるか、である。歴史的にはこれらは個別の現象を通じて段階的に見出されたが、本ノートでは以後の議論を見通しよくするため、まず公理的枠組みとしてまとめて提示する。
ここで最初に「状態」と「物理量」を分けておくことが重要である。量子論では、系がどのように準備されたかを表す情報と、測定装置が何を問うかを表す情報は同じものではない。実際、同じ準備をした系でも、位置を測るか運動量を測るか、あるいはスピンをどの軸で測るかによって、得られる結果の統計分布は変わる。

そのことを典型的に示すのがシュテルン--ゲルラッハ実験である。

\begin{experimentalfact}{シュテルン--ゲルラッハ実験}
\factitem{操作}{中性の銀原子の細いビームを、たとえば $z$ 方向に勾配をもつ不均一磁場へ通し、スクリーン上の到達位置を観測する。}
\factitem{前提}{ビームの曲がりは、磁気モーメントの $z$ 成分 $\mu_z$ に比例した力
\[
F_z \simeq \mu_z \frac{\partial B_z}{\partial z}
\]
によって生じると読む。}
\factitem{結果}{スクリーン上の像は連続的な帯ではなく、上下二つの離散的な斑点に分かれる。さらに、上側のビームだけを取り出して再び同じ $z$ 軸方向の装置に通すと分裂は起こらないが、今度は磁場の向きを $x$ 軸方向に変えた装置に通すと再び二つに分かれる。}
\factitem{示唆}{同じ準備をした系でも、どの軸の成分を測るかによって、得られる結果の統計が変わる。}
\end{experimentalfact}

したがって、測定は「もともとあった値をそのまま読む」だけではなく、「どの軸のスピン成分を問うか」という物理量そのものを指定する操作になっている。この事情を最も簡潔に表すのが、「状態はベクトルで、物理量はその上に作用する演算子で表す」という量子力学の公理的定式化である。

\subsection{状態（公理I）}
系の物理的状態は、複素ヒルベルト空間の元（状態ベクトル）$\psi$ によって表される。ヒルベルト空間とは、内積 $\langle \psi,\varphi\rangle$ が定義された複素ベクトル空間であって、その内積から定まるノルムに関して完備なものである。量子論で複素数を用いるのは、位相差が干渉を支配するからである。$\psi$ が複素係数倍（全体位相）の違いを除いて同一なら、物理的に同一の状態である。ノルムが1になるような係数を選ぶことを「正規化」といい、
\[
\langle \psi,\psi\rangle =1
\]
を満たすようにする。

\subsection{観測量（公理II）}
物理量はヒルベルト空間上の自己共役演算子 $\hat{A}$ によって表される。本ノートでは、演算子であることを明示するためにハット（$\hat{\quad}$）を付す。自己共役であることは、測定値が実数になることに対応している。

\subsection{測定と確率（公理III）}
状態 $\psi$（ただし正規化されているとする）において物理量 $\hat{A}$ を測定したとき、得られる値は $\hat{A}$ の固有値のいずれかであり、その確率は状態ベクトルと固有状態の内積の絶対値の2乗（ボルンの規則）によって与えられる。

\subsection{時間発展（公理IV）}
時間発展は線形であり、かつ確率（内積）を保存するユニタリ変換として与えられる。これは、全確率
\[
\langle \psi(t),\psi(t)\rangle
\]
が時間に依らないことを意味する。具体的なヒルベルト空間の選び方は、各理論を導入する段階でその都度与える。
